\chapter{Conclusion}

Le module \code{auth-core} (YowAuth) fournit un système d'authentification robuste, moderne et standardisé pour l'écosystème RT-Comops. En appliquant les principes de l'architecture hexagonale et en tirant profit de Spring Boot WebFlux, il garantit un traitement hautement performant et découplé des problématiques de transport et de persistance.

Les travaux menés au cours de ce projet à l'École Nationale Supérieure Polytechnique de Yaoundé (ENSPY) sous la direction du Pr. Eng Thomas Djotio et la coordination de M. Tsafack Fotso Savio ont permis d'atteindre l'ensemble des objectifs fixés :
\begin{itemize}
  \item Une modélisation stricte de l'agrégat \code{UserAccount} découplé des entités physiques de profil.
  \item La mise en œuvre d'un parcours de connexion en 2 étapes sécurisé avec détection dynamique de compte.
  \item Le durcissement de la sécurité par l'intégration d'un double facteur d'authentification (MFA OTP).
  \item L'alignement sur les standards OIDC à travers le mécanisme de Token Exchange (RFC 8693).
  \item Le développement d'une application de démonstration Next.js interactive et visuelle qui fait office de vitrine d'intégration.
\end{itemize}

Les perspectives d'évolution de YowAuth incluent l'extension du support MFA aux clés de sécurité physiques (FIDO2/WebAuthn), l'intégration native avec un serveur SMTP pour l'envoi d'e-mails réels en environnement de production, et le développement complet du back-office d'administration pour la désactivation ou la suspension des comptes à distance par les administrateurs système.
